%% kapa: exam=1 round=23 year=2012 period=1 form=A subject=economics subject_ko=경제원론 seq=8 number=48
%% source: Q-Net 감정평가사 기출문제 (cst00309) artlSeq=5052578 file=raw/1008/5052578/2012년 제23회 감정평가사 제1차시험 기출문제(1교시 A형).pdf
%% conversion: kordoc 4.24.51 -> tools/gen1.py (auto); review=proofread-sonnet; flags=none
\begin{question}{48}
기업 A의 생산함수는 $Q=LK$이다. 노동과 자본의 가격이 각각 1원일 때, 다음 설명으로 옳지 않은 것은? (단, $Q$는 생산량, $L$은 노동, $K$는 자본이다.)
\begin{choices}
  \item 규모에 대한 수익이 체증한다.
  \item 노동의 한계생산은 체감한다.
  \item 자본의 양이 단기적으로 1로 고정되어 있는 경우 100개를 생산하는데 드는 총비용은 101원이다.
  \item 자본의 양이 단기적으로 1로 고정되어 있는 경우 단기총평균비용은 생산량이 늘어나면 하락한다.
  \item 자본의 양이 단기적으로 1로 고정되어 있는 경우 한계비용은 불변이다.
\end{choices}
\end{question}
\answer{2}
